\chapter{Experimental Evidence}
\label{chap:experimental}

The QOU framework (Chapter~\ref{chap:qou}) makes predictions that differ from
the standard $\Lambda$CDM cosmology and the Standard Model of particle physics.
This chapter catalogues observable signatures, organised by scale, and maps each
to existing data sets, planned observations, and potential new experiments.

% ── Large-Scale Observations ─────────────────────────────────────────────────

\section{Large-Scale Observations}
\label{sec:large-scale}

\subsection{CMB Anisotropy and $q$-Torsion Imprint}
\label{sec:cmb-torsion}

\begin{remark}[CMB non-uniformity from $q$-torsion]
\label{rem:cmb-q-torsion}
In the Descartes universe (Example~\ref{ex:descartes}), the negative Euler
characteristic $\chi(E) = -6$ forces a non-trivial distribution of angular
modes across the CR manifold (Remark~\ref{rem:descartes-cmb}).  The
$q$-torsion leaves (Definition~\ref{def:q-torsion-leaf}) imprint a
characteristic pattern on the CMB that deviates from the statistically
isotropic Gaussian field predicted by $\Lambda$CDM\@.  Specifically, the
$q$-deformed Reeb flow (Definition~\ref{def:q-reeb-flow}) introduces
\emph{torsion} into the temperature and polarisation multipoles, producing
preferred spiral directions in the $B$-mode spectrum.
\end{remark}

\paragraph{Existing data sets.}
\begin{itemize}
  \item \emph{Planck} 2018 full-sky temperature and polarisation maps
        (PR3/PR4).
  \item \emph{WMAP} nine-year data release.
  \item ACT (Atacama Cosmology Telescope) and SPT (South Pole Telescope)
        high-$\ell$ power spectra.
\end{itemize}

\paragraph{Planned observations.}
\begin{itemize}
  \item \emph{LiteBIRD} (JAXA, launch $\sim$2032): full-sky $B$-mode
        polarisation survey sensitive to large-angular-scale torsion
        signatures.
  \item \emph{CMB-S4}: ground-based next-generation CMB survey with
        improved sensitivity to polarisation anomalies.
  \item \emph{Simons Observatory}: intermediate-scale polarimetry bridging
        current and CMB-S4 sensitivities.
\end{itemize}

\paragraph{Potential new experiments.}
\begin{itemize}
  \item A dedicated search for spiral (chiral) patterns in existing CMB
        $B$-mode data, testing for the preferred handedness predicted by the
        Reeb orientation (Remark~\ref{rem:descartes-chirality}).
  \item Cross-correlation of CMB anomalies (e.g.\ the Cold Spot, hemispherical
        asymmetry) with the angular structure of $q$-torsion leaves.
\end{itemize}

\subsection{Gravitational Signatures and Emergent Geometry}
\label{sec:gravity-emergent}

\begin{remark}[Gravity as emergent substrate geometry]
\label{rem:gravity-emergent}
In the QOU framework, general relativity emerges as the substrate geometry of
the CR manifold (Remark~\ref{rem:gr-substrate-geometry}).  The quantum-weighted
stress-energy tensor $T_{\mu\nu}^{(q)}$ introduces $q$-dependent corrections
to the Einstein field equations.  This is consistent with recent work treating
gravity as arising from an information-theoretic context~\cite{carney2025}.
\end{remark}

\paragraph{Existing data sets.}
\begin{itemize}
  \item LIGO/Virgo/KAGRA gravitational wave catalogues (GWTC-1 through
        GWTC-3).
  \item Pulsar Timing Array data sets (NANOGrav 15-year, EPTA, PPTA,
        IPTA).
  \item Solar System ephemeris tests of general relativity (Cassini, MESSENGER).
\end{itemize}

\paragraph{Planned observations.}
\begin{itemize}
  \item \emph{LISA} (ESA, launch $\sim$2035): space-based gravitational
        wave detector sensitive to low-frequency corrections from the
        $q$-deformed field equations.
  \item \emph{Einstein Telescope}: next-generation ground-based detector
        with sensitivity to post-merger ringdown deviations.
\end{itemize}

\paragraph{Potential new experiments.}
\begin{itemize}
  \item Precision tests of the equivalence principle at quantum scales,
        probing whether the substrate parameter $q$
        (Definition~\ref{def:substrate-parameter}) introduces
        composition-dependent deviations.
  \item Tabletop gravity experiments at sub-millimetre scales to search for
        $q$-dependent modifications to the Newtonian potential.
\end{itemize}

\subsection{Large-Scale Structure and Topological Defects}
\label{sec:large-scale-structure}

\begin{remark}[Topological signatures in large-scale structure]
\label{rem:topological-defects}
The $\mathrm{SO}(3)$ holonomy of the bulk bundle
(Definition~\ref{def:c-quantum-universe}) and the monodromy fixed points of the
quantum characteristic classes (Definition~\ref{def:quantum-characteristic-classes})
may leave topological imprints on the large-scale distribution of matter --- for
instance, preferred filamentary alignments or voids correlated with the brane
tower structure (Definition~\ref{def:recursive-brane-tower}).
\end{remark}

\paragraph{Existing data sets.}
\begin{itemize}
  \item SDSS (Sloan Digital Sky Survey) galaxy catalogues and BAO
        measurements.
  \item DES (Dark Energy Survey) weak-lensing maps.
  \item \emph{eROSITA} X-ray sky survey cluster catalogues.
\end{itemize}

\paragraph{Planned observations.}
\begin{itemize}
  \item \emph{Euclid} (ESA): wide-field survey of galaxy shapes and
        redshifts for weak-lensing tomography.
  \item \emph{Vera C.\ Rubin Observatory} (LSST): deep optical survey
        enabling statistical tests of filament chirality.
  \item \emph{Roman Space Telescope} (NASA): high-redshift galaxy survey.
\end{itemize}

\paragraph{Potential new experiments.}
\begin{itemize}
  \item Statistical analysis of cosmic-web filament handedness as a probe
        of global $q$-torsion.
  \item Void-shape analysis to detect signatures of monodromy fixed points
        in the brane tower.
\end{itemize}

% ── Small-Scale Observations ─────────────────────────────────────────────────

\section{Small-Scale Observations}
\label{sec:small-scale}

\subsection{Mass Equation Predictions}
\label{sec:mass-predictions}

\begin{remark}[Theoretical mass ratios from the mass identity]
\label{rem:mass-ratio-predictions}
The mass identity (Definition~\ref{def:mass-identity}) expresses inertial mass
as a line integral of $\lVert\nabla_A q\rVert \cdot \lVert L \wedge F^+\rVert$
around a vortex filament.  For the Descartes universe with $\chi(E) = -6$, the
topological data of the bundle constrain the allowed vortex configurations and
hence the mass spectrum.  The electron--muon mass ratio, for instance, is
governed by the ratio of vertical degrees in the Reeb decomposition
(Remark~\ref{rem:descartes-leptons}).  Evaluating the mass identity at
successive levels of the brane tower
(Definition~\ref{def:recursive-brane-tower}) may yield theoretical predictions
for lepton and hadron mass ratios.
\end{remark}

\paragraph{Existing data sets.}
\begin{itemize}
  \item CODATA recommended values of fundamental constants (electron, muon,
        tau masses; fine-structure constant).
  \item Particle Data Group (PDG) Review of Particle Physics.
  \item Precision atomic spectroscopy measurements (hydrogen, muonium,
        positronium).
\end{itemize}

\paragraph{Planned observations.}
\begin{itemize}
  \item Muon $g-2$ experiment at Fermilab (E989): precision measurement of
        the anomalous magnetic moment.
  \item COMET and Mu2e experiments: searches for charged lepton flavour
        violation that would constrain inter-generation mixing in the Reeb
        decomposition.
\end{itemize}

\paragraph{Potential new experiments.}
\begin{itemize}
  \item High-precision comparisons of the electron--muon and muon--tau mass
        ratios against values computed from the mass identity with specific
        $q$ and $\chi(E)$.
  \item Lattice calculations of QOU vortex masses to compare with QCD
        lattice results.
\end{itemize}

\subsection{Quantum Hall Effect of Photons}
\label{sec:quantum-hall-photons}

\begin{remark}[Photonic quantum Hall effect and $q$-deformed braiding]
\label{rem:photonic-hall}
Recent experiments have demonstrated a quantized Hall drift in photonic Chern
insulators~\cite{chenier2026}.  In the QOU framework, bosons are purely
horizontal $q$-harmonic forms (Definition~\ref{def:boson}), and the Hall
conductance arises from the Chern number of the bulk bundle restricted to the
horizontal distribution.  The $q$-deformation modifies the quantisation
condition, predicting fractional corrections of order $1 - q$ to the Hall
conductance.
\end{remark}

\paragraph{Existing data sets.}
\begin{itemize}
  \item Frequency-encoded photonic Chern insulator measurements of
        Ch\'enier et al.~\cite{chenier2026}.
  \item Many-body interference data in kagome lattices~\cite{guo2025}.
  \item Chiral phonon experiments in rare-earth halides~\cite{luo2023}.
\end{itemize}

\paragraph{Planned observations.}
\begin{itemize}
  \item Extension of photonic Chern insulator experiments to tunable
        $q$-regimes via lattice-geometry variation.
  \item Fractional quantum Hall experiments in synthetic photonic systems
        with controllable topology.
\end{itemize}

\paragraph{Potential new experiments.}
\begin{itemize}
  \item Measurement of Hall conductance in photonic lattices at varying
        effective $q$ parameters, seeking the $1 - q$ correction predicted
        by the quantum Planck constant
        (Definition~\ref{def:quantum-planck-constant}).
  \item Cold-atom realisations of $\mathrm{SU}(2)$ gauge fields in optical
        lattices, testing the $q$-deformed braiding statistics of
        Chapter~\ref{chap:path-integrals-braiding}.
\end{itemize}

\subsection{Lepton Anomalies and Reeb-Decomposition Tests}
\label{sec:lepton-anomalies}

\begin{remark}[Electron--muon discrepancies]
\label{rem:lepton-discrepancies}
The QOU framework predicts that lepton generations are distinguished by the
vertical degree of the Reeb decomposition
(Remark~\ref{rem:descartes-leptons}).  Discrepancies between electron and muon
interactions --- such as the anomalous magnetic moment tensions and lepton
flavour universality violations observed at the $B$-factories --- may reflect
the distinct Reeb-vertical structure of each generation rather than new
heavy-particle loops.
\end{remark}

\paragraph{Existing data sets.}
\begin{itemize}
  \item Muon $g-2$ results from BNL E821 and Fermilab E989.
  \item LHCb $R_K$ and $R_{K^*}$ measurements of lepton flavour
        universality.
  \item Belle~II early data on $B$-meson decays.
\end{itemize}

\paragraph{Planned observations.}
\begin{itemize}
  \item Full Belle~II data set ($\sim$50\,ab$^{-1}$): precision tests of
        lepton universality in $\tau$ decays.
  \item HL-LHC (High-Luminosity LHC): extended reach for rare lepton
        flavour violating processes.
\end{itemize}

\paragraph{Potential new experiments.}
\begin{itemize}
  \item Precision measurements of the proton charge radius via muonic
        vs.\ electronic hydrogen to probe generation-dependent coupling to
        the substrate parameter.
  \item Electron--muon scattering experiments (e.g.\ MUonE at CERN) to
        directly test the $q$-dependent corrections to the electromagnetic
        vertex.
\end{itemize}

\subsection{Informational and Entropic Gravity Tests}
\label{sec:entropic-gravity}

\begin{remark}[Entropic forces and the substrate parameter]
\label{rem:entropic-gravity}
The substrate parameter $q$ (Definition~\ref{def:substrate-parameter}) encodes
the coupling between the Levi form and the gauge curvature.  In the
near-classical regime $q \to 1$, gravity emerges as a geometric back-reaction
(Remark~\ref{rem:gr-substrate-geometry}).  Recent work on the quantum mechanics
of entropic forces~\cite{carney2025} provides a complementary perspective: the
entropic interpretation may correspond to the thermodynamic limit of the
$q$-deformed partition function on the brane tower.
\end{remark}

\paragraph{Existing data sets.}
\begin{itemize}
  \item Cavendish-type torsion balance experiments (e.g.\ E\"ot-Wash group).
  \item Casimir force measurements at micrometre scales.
  \item Neutron interferometry (gravity resonance spectroscopy).
\end{itemize}

\paragraph{Planned observations.}
\begin{itemize}
  \item MICROSCOPE follow-up missions: improved equivalence principle
        tests in space.
  \item Atom interferometry experiments (MAGIS, AION, ZAIGA) probing
        gravitational decoherence.
\end{itemize}

\paragraph{Potential new experiments.}
\begin{itemize}
  \item Entanglement-assisted gravity experiments testing whether
        gravitational interactions preserve quantum coherence as predicted
        by the $q$-deformed field equations.
  \item High-precision measurements of Casimir forces in geometries
        designed to isolate contributions from $q$-deformed vacuum
        fluctuations.
\end{itemize}
